\documentclass[11pt,a4paper]{article}
\usepackage[utf8]{inputenc}
\usepackage[spanish]{babel}
\usepackage[T1]{fontenc}
\usepackage{amsmath}
\usepackage{geometry}
\usepackage{booktabs}
\usepackage{longtable}
\usepackage{xcolor}
\usepackage{colortbl}
\usepackage{graphicx}
\usepackage{hyperref}
\usepackage{fancyhdr}
\usepackage{enumitem}
\usepackage{tcolorbox}
\usepackage{tikz}

\geometry{margin=2.5cm, headheight=14pt}

% Colores corporativos
\definecolor{verde}{HTML}{2E7D32}
\definecolor{verdeclaro}{HTML}{81C784}
\definecolor{gris}{HTML}{424242}
\definecolor{rojo}{HTML}{E53935}
\definecolor{azul}{HTML}{1565C0}
\definecolor{naranja}{HTML}{FF8F00}

\hypersetup{
    colorlinks=true,
    linkcolor=verde,
    urlcolor=azul,
}

% Header/Footer
\pagestyle{fancy}
\fancyhf{}
\fancyhead[L]{\textcolor{verde}{\textbf{CONSERVAR PAGA}}}
\fancyhead[R]{\textcolor{gris}{\small Documentación Técnica v1.0}}
\fancyfoot[C]{\thepage}
\fancyfoot[R]{\textcolor{gris}{\small Confidencial}}

\begin{document}

% ==================== PORTADA ====================
\begin{titlepage}
\centering
\vspace*{3cm}

{\Huge\bfseries\textcolor{verde}{CONSERVAR PAGA}}\\[0.5cm]
{\LARGE\textcolor{gris}{Dashboard de Control Financiero}}\\[2cm]

{\Large\bfseries Guía de Estructura de Archivos Excel}\\[0.3cm]
{\large Documentación Técnica para la Carga de Datos}\\[3cm]

\begin{tcolorbox}[colback=verde!5, colframe=verde, width=12cm]
\centering
\textbf{Versión:} 1.0 \quad | \quad \textbf{Fecha:} Abril 2026\\[0.2cm]
\textbf{Proyecto:} Control Financiero — Proyecto Ambiental\\
\textbf{Cliente:} Francia Ester Patiño\\
\textbf{Desarrollador:} Daniel Enrique Perdomo
\end{tcolorbox}

\vfill
\textcolor{gris}{\small Este documento es confidencial y de uso exclusivo del equipo de gestión financiera.}
\end{titlepage}

% ==================== ÍNDICE ====================
\tableofcontents
\newpage

% ==================== SECCIÓN 1 ====================
\section{Introducción}

El Dashboard de Control Financiero \textbf{CONSERVAR PAGA} procesa archivos Excel con estructura específica para generar indicadores financieros, alertas y reportes en tiempo real. Este documento describe la estructura exacta que debe tener el archivo Excel para que el sistema funcione correctamente.

\begin{tcolorbox}[colback=azul!5, colframe=azul, title=\textbf{Importante}]
\begin{itemize}[noitemsep]
    \item El archivo debe ser formato \texttt{.xlsx} (Excel 2007+).
    \item El tamaño máximo permitido es \textbf{50 MB}.
    \item Los datos se procesan en memoria y \textbf{nunca se almacenan} en el servidor.
    \item El nombre sugerido del archivo es: \texttt{ingresos y gastos Proyecto ambiental.xlsx}
\end{itemize}
\end{tcolorbox}

% ==================== SECCIÓN 2 ====================
\section{Resumen de Hojas Requeridas}

El archivo Excel debe contener las siguientes hojas. Cada nombre debe coincidir exactamente (se ignoran espacios al final).

\begin{center}
\renewcommand{\arraystretch}{1.4}
\begin{tabular}{|l|c|c|l|}
\hline
\rowcolor{verde!15}
\textbf{Nombre de Hoja} & \textbf{Obligatoria} & \textbf{Fila Header} & \textbf{Contenido} \\
\hline
\texttt{INF PAGOS} & \textcolor{rojo}{\textbf{Sí}} & Fila 4 & Pagos a contratistas y proveedores \\
\hline
\texttt{GASTOS CP} & No & Fila 1 & Presupuesto por categorías \\
\hline
\texttt{INGRESOS CP} & No & Fila 4 & Ingresos del proyecto \\
\hline
\texttt{INCENTIVOS} & No & Fila 4 & Pagos de incentivos por ciclo \\
\hline
\texttt{GASTOS VIAJE} & No & Fila 11 & Viáticos del equipo \\
\hline
\texttt{MOVIMIENTO DB Y CR} & No & Fila 3 & Movimientos bancarios (débito/crédito) \\
\hline
\texttt{FIDUCOLDEX} & No & Sin header & Flujo de caja fiduciaria \\
\hline
\texttt{OTROS EGRESOS} & No & Sin header & Egresos adicionales \\
\hline
\end{tabular}
\end{center}

% ==================== SECCIÓN 3 ====================
\section{Detalle por Hoja}

% --- 3.1 INF PAGOS ---
\subsection{INF PAGOS \textcolor{rojo}{(Obligatoria)}}

Esta es la hoja principal. Contiene todos los pagos realizados a contratistas y proveedores. El header comienza en la \textbf{fila 4}.

\subsubsection{Columnas Requeridas}

\begin{center}
\renewcommand{\arraystretch}{1.3}
\begin{longtable}{|p{4.5cm}|c|p{6.5cm}|}
\hline
\rowcolor{verde!15}
\textbf{Columna} & \textbf{Tipo} & \textbf{Descripción} \\
\hline
\endfirsthead
\hline
\rowcolor{verde!15}
\textbf{Columna} & \textbf{Tipo} & \textbf{Descripción} \\
\hline
\endhead
\texttt{Descripción Proveedor} & Texto & Nombre completo del contratista o razón social \\
\hline
\texttt{Valor Bruto} & Número & Monto bruto facturado (antes de retenciones) \\
\hline
\texttt{Valor Neto} & Número & Monto neto pagado. Si está vacío, se calcula automáticamente \\
\hline
\texttt{Fecha de Pago} & Fecha & Fecha del desembolso (formato \texttt{YYYY-MM-DD} o \texttt{DD/MM/YYYY}) \\
\hline
\texttt{CONCEPTO} & Texto & Clasificación del pago (ver tabla de conceptos abajo) \\
\hline
\texttt{Iva} / \texttt{IVA} & Número & Valor del IVA facturado \\
\hline
\texttt{Valor ReteFuente} & Número & Retención en la fuente (Art. 368 E.T.) \\
\hline
\texttt{Valor ReteIva} & Número & Retención del IVA (Art. 437-1 E.T.) \\
\hline
\texttt{Valor ReteIca} & Número & Retención del ICA (Ley 14/1983) \\
\hline
\texttt{Año} / \texttt{Ano} & Número & Año del pago \\
\hline
\texttt{Mes} & Número & Mes del pago (1--12) \\
\hline
\end{longtable}
\end{center}

\subsubsection{Cálculo del Valor Neto (Ley Colombiana)}

Si la columna \texttt{Valor Neto} está vacía o es cero, el sistema calcula automáticamente:

\begin{tcolorbox}[colback=naranja!5, colframe=naranja, title=\textbf{Fórmula — Estatuto Tributario (Art. 368)}]
$$\text{Valor Neto} = \text{Valor Bruto} - (\text{ReteFuente} + \text{ReteIva} + \text{ReteIca})$$
\end{tcolorbox}

\subsubsection{Valores Válidos para CONCEPTO}

El sistema clasifica automáticamente cada pago según el valor de la columna \texttt{CONCEPTO}:

\begin{center}
\renewcommand{\arraystretch}{1.3}
\begin{tabular}{|l|l|}
\hline
\rowcolor{verde!15}
\textbf{Valor en CONCEPTO} & \textbf{Categoría Asignada} \\
\hline
\texttt{HONORARIOS P.N.}, \texttt{SERVICIOS P.N.} & Contratistas P.N. \\
\hline
\texttt{SERVICIOS P.J.}, \texttt{HONORARIOS P.J.} & Contratistas P.J. \\
\hline
\texttt{COMPRAS}, \texttt{ARRENDAMIENTO} & Contratistas P.J. \\
\hline
\texttt{IVA}, \texttt{RETEFUENTE}, \texttt{GMF} & Impuestos \\
\hline
\texttt{INCENTIVOS} & Incentivos \\
\hline
\texttt{COMISION}, \texttt{COMISIONES} & Comisiones bancarias \\
\hline
\texttt{VIATICOS} & Gastos de viaje \\
\hline
Cualquier otro valor & Otros \\
\hline
\end{tabular}
\end{center}

\subsubsection{Ejemplo de Datos}

\begin{center}
\footnotesize
\begin{tabular}{|l|r|r|l|l|r|r|r|}
\hline
\rowcolor{verde!10}
\textbf{Desc. Proveedor} & \textbf{V. Bruto} & \textbf{V. Neto} & \textbf{Fecha Pago} & \textbf{CONCEPTO} & \textbf{ReteFte} & \textbf{ReteIva} & \textbf{ReteIca} \\
\hline
ROJAS PASTRANA E. & 5.000.000 & 4.450.000 & 2024-07-15 & HONORARIOS P.N. & 500.000 & 50.000 & 0 \\
\hline
EMPRESA ABC S.A.S & 12.000.000 & & 2024-07-20 & SERVICIOS P.J. & 420.000 & 228.000 & 115.200 \\
\hline
BANC. POPULAR & 150.000 & 150.000 & 2024-07-25 & COMISION & 0 & 0 & 0 \\
\hline
\end{tabular}
\end{center}

% --- 3.2 GASTOS CP ---
\subsection{GASTOS CP}

Contiene el presupuesto aprobado por categorías. El sistema busca filas con las descripciones de categoría para extraer los montos presupuestados.

\subsubsection{Formatos Soportados}

El sistema acepta dos formatos:

\begin{enumerate}[label=\textbf{\alph*)}]
    \item \textbf{Formato TOTAL}: Filas que contienen palabras como \texttt{TOTAL COSTOS DE PERSONAL}, \texttt{TOTAL INCENTIVOS}, etc.
    \item \textbf{Formato DESCRIPCION}: Filas con categorías tipo \texttt{A. PERSONAL}, \texttt{B. PAGO DE INCENTIVOS}, \texttt{C. PAGO DE DINAMIZADORES}, etc.
\end{enumerate}

\begin{center}
\renewcommand{\arraystretch}{1.3}
\begin{tabular}{|l|l|}
\hline
\rowcolor{verde!15}
\textbf{Palabra clave buscada} & \textbf{Categoría del Dashboard} \\
\hline
\texttt{TOTAL COSTOS DE PERSONAL} / \texttt{A. PERSONAL} & Personal \\
\hline
\texttt{TOTAL INCENTIVOS} / \texttt{B. PAGO DE INCENTIVOS} & Incentivos \\
\hline
\texttt{TOTAL COSTOS SERVICIOS BANCARIOS} / \texttt{D. COMISI\'{O}N} & Comisiones bancarias \\
\hline
\texttt{TOTAL DINAMIZADORES} / \texttt{C. PAGO DE DINAMIZADORES} & Dinamizadores \\
\hline
\texttt{TOTAL VIATICOS} / \texttt{E. PAGO DE VIATICOS} & Gastos de viaje \\
\hline
\texttt{TOTAL GASTOS DEL PROYECTO} & Total Proyecto \\
\hline
\texttt{F. RENDIMIENTOS FINANCIEROS} & Rendimientos financieros \\
\hline
\end{tabular}
\end{center}

% --- 3.3 INGRESOS CP ---
\subsection{INGRESOS CP}

Registra los ingresos recibidos por el proyecto (desembolsos oficiales, intereses, etc.). Header en \textbf{fila 4}.

\begin{center}
\renewcommand{\arraystretch}{1.3}
\begin{tabular}{|l|c|p{7cm}|}
\hline
\rowcolor{verde!15}
\textbf{Columna} & \textbf{Tipo} & \textbf{Descripción} \\
\hline
\texttt{Fecha} & Fecha & Fecha del ingreso \\
\hline
\texttt{CONCEPTO} & Texto & Descripción del ingreso \\
\hline
\texttt{VALOR} & Número & Monto del ingreso (ya neto) \\
\hline
\end{tabular}
\end{center}

% --- 3.4 INCENTIVOS ---
\subsection{INCENTIVOS}

Pagos de incentivos organizados por ciclos. Header en \textbf{fila 4}.

\begin{center}
\renewcommand{\arraystretch}{1.3}
\begin{tabular}{|l|c|p{7cm}|}
\hline
\rowcolor{verde!15}
\textbf{Columna} & \textbf{Tipo} & \textbf{Descripción} \\
\hline
\texttt{PAGO INCENTIVOS} / \texttt{CICLO} & Texto & Identificador del ciclo de pago \\
\hline
\texttt{FECHA DE PAGO} & Fecha & Fecha del desembolso \\
\hline
\texttt{VALOR PAGO TOTAL} & Número & Monto bruto del incentivo \\
\hline
\texttt{ABONO NETO} & Número & Monto neto recibido por el beneficiario \\
\hline
\end{tabular}
\end{center}

% --- 3.5 GASTOS VIAJE ---
\subsection{GASTOS VIAJE}

Viáticos del equipo de trabajo. Header en \textbf{fila 11}.

\begin{center}
\renewcommand{\arraystretch}{1.3}
\begin{tabular}{|l|c|p{7cm}|}
\hline
\rowcolor{verde!15}
\textbf{Columna} & \textbf{Tipo} & \textbf{Descripción} \\
\hline
\texttt{NOMBRE COMPLETO} & Texto & Nombre del funcionario \\
\hline
\texttt{PAGO VIATICOS} & Número & Monto del viático pagado \\
\hline
\texttt{FECHA DE PAGO} & Fecha & Fecha del pago \\
\hline
\end{tabular}
\end{center}

% --- 3.6 FIDUCOLDEX ---
\subsection{FIDUCOLDEX}

Flujo de caja proyectado de la fiduciaria. \textbf{Sin header} (se lee como datos crudos).

\begin{tcolorbox}[colback=verde!5, colframe=verde]
\textbf{Estructura esperada:}
\begin{itemize}[noitemsep]
    \item Fila 3: Nombres de los meses (columnas C a Q).
    \item Fila con texto \texttt{SALDO FINAL}: Saldos proyectados por mes.
\end{itemize}
\end{tcolorbox}

% --- 3.7 MOVIMIENTO DB Y CR ---
\subsection{MOVIMIENTO DB Y CR}

Movimientos bancarios (débitos y créditos). Header en \textbf{fila 3}. Se utiliza para calcular saldos mensuales.

\begin{center}
\renewcommand{\arraystretch}{1.3}
\begin{tabular}{|l|c|p{7cm}|}
\hline
\rowcolor{verde!15}
\textbf{Columna} & \textbf{Tipo} & \textbf{Descripción} \\
\hline
\texttt{FECHA} & Fecha & Fecha del movimiento bancario \\
\hline
\texttt{DEBITO} / \texttt{D\'{E}BITO} & Número & Monto debitado \\
\hline
\texttt{CREDITO} / \texttt{CR\'{E}DITO} & Número & Monto acreditado \\
\hline
\texttt{SALDO} / \texttt{SALDO FINAL} & Número & Saldo después del movimiento \\
\hline
\end{tabular}
\end{center}

% ==================== SECCIÓN 4 ====================
\section{Reglas de Normalización Automática}

El sistema aplica las siguientes transformaciones al procesar los datos:

\begin{enumerate}
    \item \textbf{Textos}: Se convierten a mayúsculas y se eliminan espacios dobles.\\
    \texttt{"ROJAS~~PASTRANA~~ESTEBAN~"} $\rightarrow$ \texttt{"ROJAS PASTRANA ESTEBAN"}
    
    \item \textbf{Fechas}: Se aceptan los formatos \texttt{YYYY-MM-DD}, \texttt{DD/MM/YYYY} y \texttt{YYYY/MM/DD}. Las filas sin fecha válida se descartan.
    
    \item \textbf{Valores numéricos}: Se convierten a tipo numérico. Los valores no numéricos se tratan como \texttt{0}.
    
    \item \textbf{Categorías}: Se normalizan automáticamente basándose en el campo \texttt{CONCEPTO} de la hoja \texttt{INF PAGOS}.
\end{enumerate}

% ==================== SECCIÓN 5 ====================
\section{Seguridad y Confidencialidad}

\begin{tcolorbox}[colback=rojo!5, colframe=rojo, title=\textbf{Política de Protección de Datos}]
\begin{itemize}[noitemsep]
    \item \textbf{Procesamiento en memoria}: Los datos del Excel se procesan exclusivamente en la memoria RAM del servidor. No se crean copias en disco.
    \item \textbf{Sin persistencia}: Al cerrar la sesión o la pestaña del navegador, todos los datos se eliminan de memoria.
    \item \textbf{Validación de archivo}: Antes de procesar, el sistema verifica la extensión (\texttt{.xlsx}), el tamaño ($\leq$ 50 MB) y la firma binaria del archivo (magic bytes).
    \item \textbf{Autenticación}: Acceso protegido con \textbf{bcrypt} + \textbf{JWT} (8 horas de validez).
    \item \textbf{Auditoría}: Cada carga de archivo se registra con hash SHA-256 en \texttt{audit\_log.jsonl}.
    \item \textbf{OWASP Top 10}: Protección contra fuerza bruta (5 intentos, 5 min lockout), timeout por inactividad (30 min), y sanitización de inputs.
\end{itemize}
\end{tcolorbox}

% ==================== SECCIÓN 6 ====================
\section{Diagrama de Flujo de Procesamiento}

\begin{center}
\begin{tcolorbox}[colback=verde!5, colframe=verde, width=12cm, title=\textbf{Pipeline de Procesamiento}]
\begin{enumerate}[noitemsep]
    \item \textbf{Subida del archivo} --- El usuario carga el archivo \texttt{.xlsx}
    \item \textbf{Validaci\'{o}n de seguridad} --- Extensi\'{o}n, tama\~{n}o y magic bytes
    \item \textbf{Carga en memoria} --- \texttt{load\_excel()} lee todas las hojas
    \item \textbf{Procesamiento} --- \texttt{process\_payments()} consolida pagos
    \item \textbf{Cubo anal\'{i}tico} --- \texttt{build\_analytics\_cube()} normaliza datos
    \item \textbf{Filtros din\'{a}micos} --- El usuario selecciona fechas, categor\'{i}as, proveedores
    \item \textbf{Renderizado} --- KPIs, gr\'{a}ficos y tablas se generan en tiempo real
\end{enumerate}

\vspace{0.3cm}
\centering
\small\textcolor{gris}{Cada paso se ejecuta en memoria. No se almacena nada en disco.}
\end{tcolorbox}
\end{center}

% ==================== SECCIÓN 7 ====================
\section{Preguntas Frecuentes}

\begin{enumerate}
    \item \textbf{¿Qué pasa si una hoja opcional no existe?}\\
    El sistema muestra una advertencia en la barra lateral pero continúa funcionando con las hojas disponibles.
    
    \item \textbf{¿Puedo usar formato \texttt{.xls} (Excel 97)?}\\
    Sí, el sistema acepta \texttt{.xls} y \texttt{.xlsx}, pero se recomienda \texttt{.xlsx} para mayor compatibilidad.
    
    \item \textbf{¿Qué pasa si las columnas tienen nombres ligeramente diferentes?}\\
    El sistema busca coincidencias parciales. Por ejemplo, tanto \texttt{Valor ReteFuente} como \texttt{ReteFuente} serán reconocidos.
    
    \item \textbf{¿Los datos quedan guardados en el servidor?}\\
    \textbf{No.} Todo se procesa en memoria (RAM) y se borra al cerrar la sesión.
    
    \item \textbf{¿Puedo agregar más hojas al Excel?}\\
    Sí. Las hojas no reconocidas simplemente se ignoran sin afectar el procesamiento.
\end{enumerate}

% ==================== PIE ====================
\vfill
\begin{center}
\textcolor{gris}{\rule{0.5\textwidth}{0.5pt}}\\[0.3cm]
\textcolor{verde}{\textbf{CONSERVAR PAGA}} — Dashboard de Control Financiero\\
\textcolor{gris}{\small Documento generado automáticamente · Abril 2026}
\end{center}

\end{document}
